\section{Problem 1.}

\begin{figure}[h]
    \centering
    \begin{minipage}[t]{0.45\textwidth}
        \vspace{0pt}
        A semicircle is inscribed in a quarter circle,
        and one of its diameter endpoints is a distance 1
        from the center of the quarter circle, drawn below. 
        What fraction of the quarter circle's area is contained 
        in the semicircle?  (i.e. what fraction of the quarter circle is shaded?)
    \end{minipage}
    \begin{minipage}[t]{0.45\textwidth}
        \vspace{0pt}
        \centering
        \includegraphics[width=0.5\linewidth]{p1.png}
        \caption{Given diagram.}
    \end{minipage}
\end{figure}



\begin{solution}
Something to notice is that when drawing from center of a circle to their end
results in a tangent line to the curve. This means when we draw a line from the 
bigger, outer circle the tangent point, it will also hit the center of the inner circle.
With this, we can calculate the areas symbolically. Since, we're only concerned
about the ratio or fraction, we don't actually need the exact value.

\begin{figure}[h]
    \centering
    \begin{minipage}{0.45\textwidth}
        \vspace{0pt}
        \centering
        \includegraphics[width=0.5\linewidth]{2rp1.jpg}
    \end{minipage}
    \begin{minipage}[h]{0.45\textwidth}
        \vspace{0pt}
        \begin{align*}
            A_{\text{big quarter}} &= \frac{\pi (\cancel{2}r)^2}{\cancel{4}} = \pi r^2 \\
            A_{\text{small semi}} &= \frac{\pi r^{2}}{2}
        \end{align*}
    \end{minipage}
\end{figure}


Clearly, we can see that half of the big quarter circle is enclosed by the small
semi circle.
\begin{equation*}
    \boxed{\frac{A_{\text{small semi}}}{A_{\text{big quarter}}} = \frac{1}{2}}
\end{equation*}

\end{solution}